% Source: https://artofproblemsolving.com/community/c3h1672301_distance_between_2_distinct_consecutive_points_on_a_line
% Stable identity: aops-topic-1672301-q01
\textbf{(a)} Put coordinates on the line, so that $A$ and $Z$ correspond to real numbers $a\ne z$. The midpoint $m=\dfrac{a+z}{2}$ satisfies $\min(a,z)<m<\max(a,z)$, so it lies strictly between $A$ and $Z$. Hence $A$ and $Z$ are not consecutive. Equivalently, there is no smallest positive distance between two points: if $d>0$ is a distance, then $\dfrac {d}{2}$ is a smaller one. $\blacksquare$

\textbf{(b)} The halving produces points at distances $\dfrac{1}{2^n}$ from $A$ for $n=1,2,3,\dots$; each of them is at a positive distance from $A$, and there is no ``last'' one, because after the point at distance $\dfrac{1}{2^n}$ comes the point at distance $\dfrac{1}{2^{n+1}}$. The infinite sum $\dfrac{1}{2}+\dfrac{1}{4}+\cdots=1$ describes the \emph{limit} of the distances from $B$, and that limit point is $A$ itself, not a new point $Z\ne A$. So the computation $AZ=0$ is correct for the limit, which is $A$; the error is assuming that the process ends at a point $Z$ different from $A$. Together with (a), this shows that the question ``what is the distance between two consecutive points'' has no answer: such points do not exist.
